\src{107}
\[
  \pi:(R,p_2)\times_X(R,p_1)\longrightarrow R
\]
được gọi là \emph{cấu xạ hợp thành}. Ba dữ liệu này đủ để xác định cấu trúc
phạm trù trên $X,R$; tiên đề đặt lên các dữ liệu ấy là như sau: với mọi $T$,
ba dữ liệu tương ứng cho $X(T),R(T)$ xác định trên cặp tập hợp này một cấu
trúc phạm trù (tương ứng, một cấu trúc groupoid). Khi thích hợp, điều này có
thể được diễn đạt bằng tính giao hoán của một số biểu đồ có chứa một cấu xạ
được xác định duy nhất
\[
  D:X\longrightarrow R
\]
và, trong trường hợp các groupoid, một cấu xạ được xác định duy nhất
\[
  s:R\longrightarrow R;
\]
các biểu đồ được viết tường minh đúng như trong trường hợp «theo nghĩa tập
hợp». May thay, cách diễn giải dài dòng này của các tiên đề là vô ích trong
thực hành. Lợi ích lý thuyết duy nhất của khả năng biểu diễn dữ liệu và tiên
đề bằng các cấu xạ giữa một số tích thớ và các đẳng thức giữa những cấu xạ ấy
là như sau: nếu $F:\mathcal C\to\mathcal C'$ là một hàm tử khớp trái (tức giao
hoán với các tích hữu hạn và các tích thớ), thì nó biến một
$\mathcal C$-phạm trù (tương ứng, một $\mathcal C$-groupoid) thành một
$\mathcal C'$-phạm trù (tương ứng, một $\mathcal C'$-groupoid), với điều kiện
các tích hữu hạn và các tích thớ tồn tại trong $\mathcal C$.

Trong thực hành, điều quan trọng là biết diễn giải các cấu xạ
$p_1,p_2,\pi,D,s$ như những \emph{phép toán đơn hình} trong một vật bán đơn
hình hoặc đơn hình thích hợp của $\mathcal C$, ít nhất khi các tích thớ tồn
tại trong $\mathcal C$. Để cố định thuật ngữ, ta đưa vào \emph{phạm trù
$\mathcal S$ các đơn hình chuẩn} là phạm trù mà các vật là những tập hữu hạn
có dạng
\[
  \Delta_n=[0,n],\qquad\text{trong đó }n\geq 0
\]
(khoảng các số nguyên từ $0$ đến $n$), và các cấu xạ hay mũi tên là
\emph{các ánh xạ tùy ý} giữa những tập hữu hạn ấy. Lưu ý rằng phạm trù
$\mathcal S$ tương đương với phạm trù các tập hữu hạn \emph{khác rỗng}, trong
đó các cấu xạ là những ánh xạ giữa các tập hữu hạn. Trong $\mathcal S$, đối
tích của một họ hữu hạn \emph{khác rỗng} các vật hiển nhiên tồn tại, cũng như
phép đẩy ra của hai vật dưới một vật thứ ba (phép toán đối ngẫu với tích thớ).
Ký hiệu $\mathcal S'$ là phạm trù con của $\mathcal S$ có cùng các vật, nhưng
trong đó các cấu xạ là \emph{các ánh xạ tăng} giữa các $\Delta_n$. Phạm trù
này tương đương với phạm trù các tập hữu hạn khác rỗng được sắp thứ tự toàn
phần. Trong phạm trù này,
